Gambar~\ref{fig:seq-diagram-mermaid} merupakan diagram alur proses dari sistem \textit{retrieval-augmented generation} (RAG) yang asli.

\begin{figure}[H]
\begin{mermaid}
%%{init: {
  'theme': 'base',
  'themeVariables': {
    'primaryColor': '#ffffff',
    'primaryBorderColor': '#000000',
    'primaryTextColor': '#000000',
    'lineColor': '#000000',
    'textColor': '#000000',
    'clusterBkg': '#ffffff',
    'clusterBorderColor': '#000000',
    'fontSize': '11px',
    'fontFamily': 'Times New Roman, Times, serif'
  },
  'flowchart': {
    'curve': 'linear',
    'nodeSpacing': 40,
    'rankSpacing': 40,
    'padding': 12
  }
}}%%
graph TD
    classDef formal fill:#ffffff,stroke:#000000,stroke-width:1.5px;

    NodeStart(["Mulai: Input Pertanyaan <code>x</code>"]):::formal --> ProcInput["Enkode Pertanyaan <code>x</code>"]:::formal

    subgraph Sub_Retrieval ["TAHAP 1: RETRIEVAL (DPR)"]
        ProcInput --> RetData["Cari Chunks Relevan"]:::formal
        WikiIndex[("Indeks Wikipedia:<br/>21M Chunks 100-Kata")]:::formal
        RetTopK["Ambil Top-K Chunks <code>z</code>"]:::formal
        
        WikiIndex -.->|Pencarian MIPS| RetData
        RetData --> RetTopK
    end

    RetTopK --> ProcContext

    subgraph Sub_Generation ["TAHAP 2: GENERATION (BART)"]
        ProcContext["Gabungkan Pertanyaan <code>x</code><br/>+ Chunks <code>z</code>"]:::formal --> ModelBART["Generate Jawaban <code>y</code>"]:::formal
        ParamMem[("Memori Model:<br/>BART-large 400M")]:::formal
        ParamMem -.->|Pre-trained Weights| ModelBART
    end

    ModelBART --> NodeEnd(["Selesai: Output Jawaban <code>y</code>"]):::formal
\end{mermaid}
\caption{Sequence diagram (numbering, asynchronous arrows, multiple participants)}
\label{fig:seq-diagram-mermaid}
\end{figure}